% Part: proof-theory
% Chapter: sequent-calculus
% Section: proof-examples

\documentclass[../../../include/open-logic-section]{subfiles}

\begin{document}

\olfileid{pt}{seq}{ex}

\olsection{证明举例}

要给相继式构造推导，通常方便的做法是从目标相继式开始向上回溯：
选择其中的一个公式作为活动公式，在相继式上方写出相应规则的前提，
然后重复这一过程，直至得到公理。例如，假设我们要证明相继式
\[(!C \land !D) \lif !E \Sequent !C \lif (!D \lif !E)\]
考虑 $!C \lif (!D \lif !E)$：它具有 $!A \lif !B$ 的形式，并且出现在相继式的右边。
将整个相继式与 \Log{G1c} 的 $\RightR{\lif}$ 规则进行匹配，
\[\Axiom$ !A, \Gamma \fCenter \Delta, !B$
\RightLabel{\RightR{\lif}}
\UnaryInf$ \Gamma \fCenter \Delta, !A \lif !B $
\DisplayProof\]
这提示我们应该取 $\Gamma = \{!C \lif (!D \lif !E)\}$，$\Delta =
\emptyset$，$!A \ident !C$ 和 $!B \ident !D \lif !E$。因此，证明的最后一步推理可以是
\[\Axiom$ !C, (!C \land !D) \lif !E \fCenter !D \lif !E$
\RightLabel{\RightR{\lif}}
\UnaryInf$(!C \land !D) \lif !E \fCenter !C \lif (!D \lif !E)$
\DisplayProof\]
如果我们重复这一思路，将 $!D \lif !E$ 作为活动公式，便得到
\[\Axiom$ !D, !C, (!C \land !D) \lif !E \fCenter !E$
\RightLabel{\RightR{\lif}}
\UnaryInf$!C, (!C \land !D) \lif !E \fCenter !D \lif !E$
\RightLabel{\RightR{\lif}}
\UnaryInf$(!C \land !D) \lif !E \fCenter !C \lif (!D \lif !E)$
\DisplayProof\]
现在我们需要关注 $(!C \land !D) \lif !E$，并使用 \Log{G1c} 的 $\LeftR{\lif}$ 规则，
\[\Axiom$ \Gamma \fCenter \Delta, !A$
\Axiom$ !B, \Gamma \fCenter \Delta$
\RightLabel{\LeftR{\lif}}
\BinaryInf$ !A \lif !B, \Gamma \fCenter \Delta$
\DisplayProof\]
这引导我们得到
\[\Axiom$!D, !C \fCenter !E, !C \land !D$
\Axiom$!D, !C, !E \fCenter !E$
\RightLabel{\LeftR{\lif}}
\BinaryInf$ !D, !C, (!C \land !D) \lif !E \fCenter !E$
\RightLabel{\RightR{\lif}}
\UnaryInf$!C, (!C \land !D) \lif !E \fCenter !D \lif !E$
\RightLabel{\RightR{\lif}}
\UnaryInf$(!C \land !D) \lif !E \fCenter !C \lif (!D \lif !E)$
\DisplayProof\]
将最左顶端相继式中的公式 $!C \land !D$ 作为主公式，并应用 $\RightR{\land}$ 规则，得到
\[
\Axiom$!D, !C \fCenter !E, !C$
\Axiom$!D, !C \fCenter !E, !D$
\RightLabel{\RightR{\land}}
\BinaryInf$!D, !C \fCenter !E, !C \land !D$
\Axiom$!D, !C, !E \fCenter !E$
\RightLabel{\RightR{\lif}}
\BinaryInf$ !D, !C, (!C \land !D) \lif !E \fCenter !E$
\RightLabel{\RightR{\lif}}
\UnaryInf$!C, (!C \land !D) \lif !E \fCenter !D \lif !E$
\RightLabel{\RightR{\lif}}
\UnaryInf$(!C \land !D) \lif !E \fCenter !C \lif (!D \lif !E)$
\DisplayProof
\]
目前为止一切顺利。注意，我们以上所做的在 \Log{G3c} 中同样可行，因为迄今为止我们用到的规则在 \Log{G1c}
和~\Log{G3c} 中是相同的。我们得到了一个推导片段，其中所有顶端相继式左右两边都包含同一个公式。
但它们还不算公理。

在 \Log{G1c} 中，我们必须从形如 $!A \Sequent !A$ 的公理出发，给出最顶端那些相继式的推导。
这很容易用弱化规则完成，例如，最左顶端的相继式 $!D, !C \Sequent
!E, !C$ 可以如下推导：
\[
\Axiom$!C \fCenter !C$
\RightLabel{\LeftR{\Weakening}}
\UnaryInf$!D, !C \fCenter !C$
\RightLabel{\RightR{\Weakening}}
\UnaryInf$!D, !C \fCenter !E, !C$
\]
综合来看，我们在 \Log{G1c} 中得到了如下推导：
\[
\Axiom$!C \fCenter !C$
\RightLabel{\LeftR{\Weakening}}
\UnaryInf$!D, !C \fCenter !C$
\RightLabel{\RightR{\Weakening}}
\UnaryInf$!D, !C \fCenter !E, !C$
\Axiom$!D \fCenter !D$
\RightLabel{\LeftR{\Weakening}}
\UnaryInf$!D, !C \fCenter !D$
\RightLabel{\RightR{\Weakening}}
\UnaryInf$!D, !C \fCenter !E, !D$
\RightLabel{\RightR{\land}}
\BinaryInf$!D, !C \fCenter !E, !C \land !D$
\Axiom$!E \fCenter !E$
\RightLabel{\LeftR{\Weakening}}
\UnaryInf$!C, !E \fCenter !C$
\RightLabel{\LeftR{\Weakening}}
\UnaryInf$!D, !C, !E \fCenter !E$
\RightLabel{\LeftR{\lif}}
\BinaryInf$ !D, !C, (!C \land !D) \lif !E \fCenter !E$
\RightLabel{\RightR{\lif}}
\UnaryInf$!C, (!C \land !D) \lif !E \fCenter !D \lif !E$
\RightLabel{\RightR{\lif}}
\UnaryInf$(!C \land !D) \lif !E \fCenter !C \lif (!D \lif !E)$
\DisplayProof
\]
这个证明的结构（特别是其高度）与公式 $!C$、$!D$ 和~$!E$ 的具体形式无关。
我们可以用任意公式替换上述推导中的这些符号，结果仍然是~\Log{G1c} 中的一个推导。
\Log{G1c} 中的这个证明是\emph{模式化的}。

在 \Log{G3c} 中，公理是形如 $!A, \Gamma \Sequent
\Delta, !A$ 的相继式，其中 $!A$ 必须是原子公式。因此，尽管 $!D, !C \Sequent
!E, !C$ 看起来像 \Log{G3c} 的一个公理（$\Gamma = \{!D\}$；$\Delta =
\{!E\}$），但仅当 $!C$ 实际上是原子公式时，它才\emph{真正是}一个公理。
我们的推导片段只有当 $!C$、
$!D$ 和 $!E$ 都是原子公式时，才是 \Log{G3c} 中的一个完整证明。

虽然严格来说，我们还没有证明 \[\Log{G1c} \Proves (!C \land !D) \lif !E \fCenter !C \lif (!D \lif
!E),\]，但实际上这就是我们需要（并且能够）做的全部工作。这是因为，每一个形如 $!A, \Gamma \Sequent \Delta, !A$ 的相继式在 \Log{G3c} 中都是可证的，即使 $!A$ 本身不是原子公式。我们可以通过对~$!A$ 的深度~$\depth{!A}$ 进行归纳来证明这一点。

\begin{defn}\ollabel{defn:depth}
公式的\emph{深度}~$\depth{!A}$ 定义如下：
\begin{enumerate}
  \item 若 $!A$ 是原子公式，则 $\depth{!A} = 0$。
  \item 若 $!A \ident \lnot !B$、$!A \ident \lforall[x][A]$ 或
  $!A \ident \lexists[x][A]$，则 $\depth{!A} = \depth{!B} + 1$。
  \item 若 $!A \ident (!B \land !C)$、$!A \ident (!B \lor !C)$ 或 $!A \ident (!B \lif
  !C)$，则 $\depth{!A} = \max(\depth{!B},\depth{!C}) + 1$。
\end{enumerate}
\end{defn}

不难证明，对于任意项~$t$，$\depth{A(x)} = \depth{A(t)}$。关键在于，在任何规则中，主公式的深度总是大于活动公式的深度。例如，我们有：
\begin{align*}
\depth{P(x)} = \depth{P(a)} & = 0,\\
\depth{\lforall[x][P(x)]} & = 1,\\
\depth{(\lforall[x][P(x)] \lif P(a))} & = \max(1,0) + 1 = 2.
\end{align*}
我们可以利用这一点，通过对~$\depth{!A}$ 进行归纳来证明一些结果，如下命题。

\begin{prop}\ollabel{prop:G1c-axioms}
对于任意~$!A$，在 $\Log{G1c}$ 中存在一个 $!A \Sequent !A$ 的推导，其中所有公理都是原子公式。
\end{prop}

\begin{proof}
  对~$\depth{!A}$ 进行归纳。若 $\depth{!A} = 0$，则 $!A$ 是原子公式，
  此时 $!A \Sequent !A$ 已经是~\Log{G1c} 中符合要求形式的推导。
  若 $\depth{!A} > 0$，则 $!A$ 由深度较低的公式构造而成，
  例如 $!A \ident \lnot !B$、$!A \ident
  (!B \land !C)$，或 $!A \ident \lforall[x][A(x)]$。
  归纳假设为：对所有满足 $\depth{!D} < \depth{!A}$ 的 $!D$，
  在 $\Log{G1c}$ 中都能从原子公理推出 $!D \Sequent !D$。

  若 $!A \ident \lnot !B$，则 $\depth{!B} < \depth{!A}$，
  根据归纳假设，在 \Log{G1c} 中存在从原子公理出发的
  $!B \Sequent !B$ 推导~$\pi_1$。
  我们可以如下将其扩展为 $\lnot !B \Sequent \lnot !B$ 的推导：
  \[
  \AxiomC{}
  \RightLabel{$\pi_1$}
  \Deduce$!B \fCenter !B$
  \RightLabel{\LeftR{\lnot}}
  \UnaryInf$\lnot !B, !B \fCenter $
  \RightLabel{\RightR{\lnot}}
  \UnaryInf$\lnot !B \fCenter \lnot !B$
  \DisplayProof
\]

  若 $!A \ident (!B \land !C)$，则 $\depth{!B} < \depth{!A}$ 且
  $\depth{!C} < \depth{!A}$。根据归纳假设，在 \Log{G1c} 中存在从原子公理出发的
  $!B \Sequent !B$ 和 $!C \Sequent !C$ 的推导~$\pi_1$ 和~$\pi_2$。
  我们可以如下将它们扩展为 $!B \land !C \Sequent !B \land !C$ 的推导：
  \[
  \AxiomC{}
  \RightLabel{$\pi_1$}
  \Deduce$!B, \fCenter !B$
  \RightLabel{\LeftR{\Weakening}}
  \UnaryInf$!B, !C, \fCenter !B$
  \RightLabel{\LeftR{\land}}
  \UnaryInf$!B \land !C, !C\fCenter !B$
  \RightLabel{\LeftR{\land}}
  \UnaryInf$!B \land !C, !B \land !C \fCenter !B$
  \AxiomC{}
  \RightLabel{$\pi_2$}
  \Deduce$!C, \fCenter !C$
  \RightLabel{\LeftR{\Weakening}}
  \UnaryInf$!B, !C, \fCenter !C$
  \RightLabel{\LeftR{\land}}
  \UnaryInf$!B \land !C, !C \fCenter !C$
  \RightLabel{\LeftR{\land}}
  \UnaryInf$!B \land !C, !B \land !C \fCenter !C$
  \RightLabel{\RightR{\land}}
  \BinaryInf$!B \land !C \fCenter !B \land !C$
  \RightLabel{\LeftR{\Contraction}}
  \UnaryInf$!B \land !C \fCenter !B$
  \DisplayProof
  \]

  现在假设 $!A \ident \lforall[x][B(x)]$。取一个不在 $\lforall[x][B(x)]$ 中出现的常项符号~$c$。
  则 $\depth{B(c)} =
  \depth{!B(x)} < \depth{!A}$，根据归纳假设，在 \Log{G1c} 中存在从原子公理出发的
  $!B(c)
  \Sequent !B(c)$ 推导~$\pi_1$。
  我们可以如下将其扩展为 $\lforall[x][B(x)] \Sequent \lforall[x][B(x)]$ 的推导：
  \[
  \AxiomC{}
  \RightLabel{$\pi_1$}
  \Deduce$!B(c) \fCenter !B(c)$
  \RightLabel{\LeftR{\lforall}}
  \UnaryInf$\lforall[x][!B(x)] \fCenter !B(c)$
  \RightLabel{\RightR{\lforall}}
  \UnaryInf$\lforall[x][!B(x)] \fCenter \lforall[x][!B(x)]$
  \DisplayProof
\]

  其余情况留作练习。
\end{proof}

\begin{prob}
  通过处理 $!A \ident (!B \lor !C)$、$!A \ident (!B \lif !C)$ 以及
  $!A \ident \lexists[x][B(x)]$ 的情况，补全 \olref{prop:G1c-axioms} 的证明。
\end{prob}

就我们的目的而言，在 $!A \ident \lnot !B$ 的情形下，本也可以使用推导
\[
\AxiomC{}
\RightLabel{$\pi_1$}
\Deduce$!B \fCenter !B$
\RightLabel{\RightR{\lnot}}
\UnaryInf$\fCenter !B, \lnot !B$
\RightLabel{\LeftR{\lnot}}
\UnaryInf$\lnot !B \fCenter \lnot !B$
\DisplayProof
\]。
然而，在某些相继式系统中，这样做是行不通的。
直觉主义相继式演算 \Log{G1i} 类似于 \Log{G1c}，只是它限制所有相继式的右边至多有一个公式。
若 $\Delta = \emptyset$，则第一种推导也可以是~\Log{G1i} 中的推导，
但第二种则不行，因为其中一个相继式的右边同时包含 $!B$ 和~$\lnot !B$。

请注意，即使在 \Log{G1c} 中，对于 $!A \ident \lforall[x][!B(x)]$ 的情形，我们也无法采用不同的规则顺序。
以下\emph{不是}推导：
  \[
  \AxiomC{}
  \RightLabel{$\pi_1$}
  \Deduce$!B(c) \fCenter !B(c)$
  \RightLabel{\RightR{\lforall}}
  \UnaryInf$!B(c) \fCenter \lforall[x][!B(x)]$
  \RightLabel{\LeftR{\lforall}}
  \UnaryInf$\lforall[x][!B(x)] \fCenter \lforall[x][!B(x)]$
  \DisplayProof
\]
因为它违反了 $\RightR{\lforall}$ 规则的特征变元条件。

\begin{prop}\ollabel{prop:G3c-axioms}
任何形如 $!A, \Gamma
\Sequent \Delta, !A$（其中 $!A$ \emph{未必}是原子公式）的相继式
在~\Log{G3c} 中都是可证的。
\end{prop}

\begin{proof}
证明与 \olref{prop:G1c-axioms} 的证明非常相似，
但在归纳假设上必须小心。当 $n =
\depth{!A} = 0$ 时的情形同样简单：若 $n=0$，则 $!A$ 是原子公式，
于是 $!A, \Gamma \Sequent \Delta, !A$ 是~\Log{G3c} 的公理。

现在假设 $n>0$。我们再次分情况讨论。归纳假设是：
对所有满足 $\depth{!D} < n$ 的 $!D$，以及所有的 $\Pi$ 和 $\Lambda$，
$\Log{G3c} \Proves !D, \Pi \Sequent \Lambda, !D$。
以下是 $!A \ident (!B \land !C)$ 的情形：我们两次使用归纳假设：
一次取 $!D = !B$、$\Pi = !C, \Gamma$ 和 $\Lambda
= \Delta$，得到 $!B, !C, \Gamma \Sequent
\Delta, !B$ 的推导~$\pi_1$；
另一次取 $!D = !C$、$\Pi = !B, \Lambda$ 和 $\Lambda
= \Delta$，得到 $!B, !C, \Gamma \Sequent
\Delta, !C$ 的推导~$\pi_2$。
以下是 $!B \land !C, \Gamma \Sequent
\Delta, !B \land !C$ 的推导：
\[
\AxiomC{}
\RightLabel{$\pi_1$}
\Deduce$!B, !C, \Gamma \fCenter \Delta, !B$
\RightLabel{\LeftR{\land}}
\UnaryInf$!B \land !C, \Gamma \fCenter \Delta, !B$
\AxiomC{}
\RightLabel{$\pi_2$}
\Deduce$!B, !C, \Gamma \fCenter \Delta, !C$
\RightLabel{\LeftR{\land}}
\UnaryInf$!B \land !C, \Gamma \fCenter \Delta, !C$
\RightLabel{\RightR{\land}}
\BinaryInf$!B \land !C, \Gamma \fCenter \Delta, !B \land !C$
\DisplayProof
\]

对于 $!A \ident \lforall[x][!B(x)]$ 的情形，选取一个不在
$\lforall[x][!B(x)]$、$\Gamma$ 或~$\Delta$ 中出现的常项符号。
对 $!D = !B(c)$、$\Pi = \lforall[x][!B(x)], \Gamma$
以及 $\Lambda = \Delta$ 应用归纳假设。
由此得到 $!B(c),
\lforall[x][!B(x)], \Gamma \Sequent \Delta, !B(c)$ 的推导~$\pi_1$，据此我们得到：
\[
\AxiomC{}
\RightLabel{$\pi_1$}
\Deduce$!B(c), \lforall[x][!B(x)], \Gamma \fCenter \Delta, !B(c)$
\RightLabel{\LeftR{\lforall}}
\UnaryInf$\lforall[x][!B(x)], \Gamma \fCenter \Delta, !B(c)$
\RightLabel{\RightR{\lforall}}
\UnaryInf$\lforall[x][!B(x)], \Gamma \fCenter \Delta, \lforall[x][!B(x)]$
\DisplayProof
\]
我们选取~$c$ 的方式使得 $\RightR{\lforall}$ 规则的特征变元条件得以满足。
\end{proof}

\begin{prob}
  通过处理 $!A \ident (!B \lif !C)$ 和 $!A \ident
  \lexists[x][B(x)]$ 的情形，
  继续完成 \olref{prop:G3c-axioms} 的证明。
\end{prob}

\end{document}